\input{header.tex}
\input{newcommands.tex}


\title{Stoffverteilungspläne automatisch erstellen}
\author{Fortbildungsmaterial für Lehrkräfte}
\date{\today}
\pagestyle{plain}
\setlength{\parskip}{1.5em} % Abstand zwischen Absätzen (Handreichungen)
\setlength{\emergencystretch}{4em} % lange \texttt-Wörter umbrechen statt über den Rand


\begin{document}

\maketitle

Ein Stoffverteilungsplan beantwortet eine einfache Frage mit viel Rechenarbeit: Welches Thema behandle ich an welchem Tag in welcher Klasse? Dazwischen liegen Ferien, Feiertage, A- und B-Wochen, Klassenfahrten und Prüfungszeiträume. Das Python-Skript \texttt{stoffverteilungsplan.py} übernimmt diese Arbeit. Man beschreibt Klassen, Termine, Sperren und Lernbereiche in einer kurzen Textdatei; das Skript verteilt die Lernbereiche \textbf{strikt nach Unterrichtsstunden} auf die echten Termine jeder Klasse und schreibt einen fertigen \hologo{LaTeX}-Plan. Im Terminal meldet es, ob der Pflichtumfang bis zum Schuljahresende in jede Klasse passt.

\tableofcontents



%%%%  KAPITEL 1  %%%%
\section{Das Prinzip}

\begin{center}
\begin{tabular}{@{}>{\raggedright\arraybackslash}p{4.3cm}>{\raggedright\arraybackslash}p{5.3cm}>{\raggedright\arraybackslash}p{5.4cm}@{}}
	\toprule
	\textbf{Du änderst} & \textbf{Werkzeug} & \textbf{Ergebnis} \\
	\midrule
	\texttt{plan\_\dots.toml} (Klassen, Termine, Lernbereiche) & \texttt{stoffverteilungsplan.py} & \texttt{plan\_\dots.tex}, daraus mit \texttt{pdflatex} der Plan \\
	\texttt{ferien\_\dots.toml} (einmal pro Schuljahr) & & \\
	\bottomrule
\end{tabular}
\end{center}

Man muss den Code des Skripts nicht verstehen: Eingabedatei ändern, Skript starten, PDF übersetzen. Gestartet wird im Terminal im Ordner der Dateien:

\begin{lstlisting}[style=latex]
python3 stoffverteilungsplan.py plan_informatik10.toml
pdflatex plan_informatik10.tex
\end{lstlisting}

\textbf{Einmalig einrichten:} Python ab Version~3.11 (wegen des eingebauten TOML-Lesers) und eine \hologo{LaTeX}-Installation. Weitere Python-Pakete sind nicht nötig.

\begin{hinweis}[TOML]
	Die Eingabedateien sind \textit{TOML}-Dateien: reiner Text, mit jedem Editor zu öffnen. \texttt{\#} leitet einen Kommentar ein, Texte stehen in Anführungszeichen, Datumsangaben schreibt man \texttt{2026-08-17}. Ein Block, der mehrfach vorkommen darf, steht in doppelten eckigen Klammern: \texttt{[[gruppe]]}.
\end{hinweis}



%%%%  KAPITEL 2  %%%%
\section{Die Eingabedatei}


%%%%  UNTERKAPITEL 2.1  %%%%
\subsection{Allgemeine Angaben}

\begin{lstlisting}[style=latex]
[plan]
titel      = "Stoffverteilungsplan Informatik Klasse 10"
schuljahr  = "2026/27"
lehrkraft  = "N. N."
schule     = "Musterschule"
beginn     = 2026-08-17        # erster Unterrichtstag, der geplant wird
ende       = 2027-07-09        # letzter Schultag vor den Sommerferien
uebergabe_am = 2027-02-05      # bis hier plane ich selbst (Linie im Plan); weglassen = keine Übergabe
ferien     = "ferien_sachsen_2026_27.toml"
woche_bei_beginn = "A"         # welche Woche ist die Woche von "beginn"?
ab_zaehlung = "schulwochen"    # Ferienwochen zählen nicht mit; sonst "kalenderwochen"
ustd_pro_termin = 2            # Doppelstunde
uebergabe  = "Übergabe an den Fachkollegen"
\end{lstlisting}

Mit \texttt{uebergabe\_am} zieht das Skript eine rote Linie in den Plan, etwa wenn eine Referendarin oder ein Referendar die Klasse zum Halbjahr abgibt.


%%%%  UNTERKAPITEL 2.2  %%%%
\subsection{Klassen und Termine}

Jede Klasse oder Halbgruppe ist ein Block \texttt{[[gruppe]]}. Ein Termin hat einen Wochentag, eine Woche (\texttt{"A"}, \texttt{"B"} oder \texttt{"AB"} für jede Woche) und eine Uhrzeit:

\begin{lstlisting}[style=latex]
[[gruppe]]
name    = "10b"
termine = [ { tag = "Di", woche = "A", zeit = "08:00–09:30" } ]

[[gruppe]]
name    = "10c"
termine = [ { tag = "Di", woche = "B", zeit = "08:00–09:30" } ]
\end{lstlisting}


%%%%  UNTERKAPITEL 2.3  %%%%
\subsection{Gesperrte Zeiträume}

\begin{lstlisting}[style=latex]
[[sperre]]
grund = "Klassenfahrt"
von   = 2026-09-28
bis   = 2026-10-02
modus = "ausfall"            # Termin fällt weg

[[sperre]]
grund = "Prüfungszeitraum"
von   = 2026-10-26
bis   = 2026-12-18
modus = "puffer"             # Termin bleibt, bekommt aber keinen neuen Stoff
thema = "Übung und Vertiefung (Ausfall möglich)"
\end{lstlisting}

Mit \texttt{nur = ["10b"]} gilt eine Sperre nur für einzelne Gruppen.


%%%%  UNTERKAPITEL 2.4  %%%%
\subsection{Lernbereiche}

Die Lernbereiche stehen in der Reihenfolge des Unterrichts. \texttt{ustd} sind die Unterrichtsstunden laut Lehrplan; bei Doppelstunden ergeben 12~Ustd.\ sechs Termine. Jeder Eintrag in \texttt{themen} ist ein Termin. Sind es weniger Themen als Termine, füllt das Skript mit \enquote{Fortsetzung} auf. \texttt{\{ thema = "\dots", note = true \}} markiert eine sonstige Note, \texttt{wahl = true} einen Wahlbereich, der den Rest des Schuljahres auffüllt.

\begin{lstlisting}[style=latex]
[[lernbereich]]
kurz  = "LB 1"
name  = "Algorithmen"
ustd  = 12
farbe = "akzentblau"
themen = [
  "Daten im Programm",
  "Datentypen",
  "Auswahl",
  "Wiederholung",
  { thema = "Funktionen", note = true },
  "Maschinelles Lernen",
]

[[lernbereich]]
kurz  = "WB"
name  = "Wahlbereich"
wahl  = true
farbe = "akzentlila"
themen = [ "Wahlbereich" ]
\end{lstlisting}

Die Farben sind die Akzentfarben der universalen \texttt{newcommands.tex} (\texttt{akzentblau}, \texttt{akzentgruen}, \texttt{akzentorange}, \texttt{akzentlila}, \texttt{akzentrot}).


%%%%  UNTERKAPITEL 2.5  %%%%
\subsection{Ferien und Feiertage}

Die Ferien stehen in einer eigenen Datei, die man einmal pro Schuljahr anlegt und an das ganze Kollegium weitergeben kann. Ein Eintrag ohne \texttt{bis} ist ein einzelner Tag:

\begin{lstlisting}[style=latex]
[[frei]]
name = "Herbstferien"
von  = 2026-10-12
bis  = 2026-10-24

[[frei]]
name = "Buß- und Bettag"
von  = 2026-11-18
\end{lstlisting}



%%%%  KAPITEL 3  %%%%
\section{Das Ergebnis}

Das Skript schreibt \texttt{plan\_informatik10.tex}. Nach dem Übersetzen enthält der Plan drei Teile: \textit{Teil A} mit der allgemeinen Aufteilung der Lernbereiche, eine \textit{Halbjahresübersicht} als Zeitleiste und \textit{Teil B} mit jeder Klasse, jedem Datum und der Übergabelinie. Im Terminal steht zusätzlich eine Bilanz:

\begin{lstlisting}[style=latex]
geschrieben: plan_informatik10.tex
  10b: 16 Termine mit Stoff, 4 Puffer, 0 Pflicht-Termine offen
  10d: 16 Termine mit Stoff, 4 Puffer, 0 Pflicht-Termine offen
  10c: 15 Termine mit Stoff, 4 Puffer, 0 Pflicht-Termine offen
  10e: 15 Termine mit Stoff, 4 Puffer, 0 Pflicht-Termine offen
\end{lstlisting}

\begin{ausgabe}
	\centering
	\fbox{\includegraphics[page=1, width=0.47\linewidth]{plan_informatik10.pdf}}\hfill
	\fbox{\includegraphics[page=3, width=0.47\linewidth]{plan_informatik10.pdf}}
\end{ausgabe}

\enquote{0 Pflicht-Termine offen} heißt: Alle Lernbereiche passen in jede Klasse. Passt etwas nicht, steht im Plan ein roter Hinweis, was bis zum Schuljahresende nicht untergebracht ist.



%%%%  KAPITEL 4  %%%%
\section{Was das Skript rechnet}

\begin{itemize}
	\item \textbf{Termine:} Für jede Gruppe alle Tage zwischen \texttt{beginn} und \texttt{ende}, an denen sie Unterricht hat, ohne Ferien, Feiertage und Sperren mit \texttt{modus = "ausfall"}.
	\item \textbf{A- und B-Wochen:} Ab \texttt{woche\_bei\_beginn} wechselt die Woche jede Schulwoche. Mit \texttt{ab\_zaehlung = "kalenderwochen"} wechselt sie auch in den Ferien.
	\item \textbf{Verteilung nach Stunden:} Jeder Lernbereich bekommt genau so viele Termine, wie seine Unterrichtsstunden ergeben (\texttt{ustd / ustd\_pro\_termin}). Puffer-Termine zählen nicht mit.
	\item \textbf{Wahlbereich:} Was nach dem Pflichtumfang übrig bleibt, füllt der Bereich mit \texttt{wahl = true}.
\end{itemize}



%%%%  KAPITEL 5  %%%%
\section{Zweites Beispiel: wöchentliche und 14-tägliche Stunden}

\texttt{plan\_mathe\_lk11.toml} plant einen Kurs mit zwei wöchentlichen Doppelstunden und einer zusätzlichen Doppelstunde nur in B-Wochen:

\begin{lstlisting}[style=latex]
[[gruppe]]
name    = "LK 11"
termine = [
  { tag = "Mo", woche = "AB", zeit = "10:00–11:30" },
  { tag = "Di", woche = "AB", zeit = "10:00–11:30" },
  { tag = "Fr", woche = "B",  zeit = "12:30–14:00", bemerkung = "14-täglich" },
]
\end{lstlisting}

\begin{ausgabe}
	\centering
	\fbox{\includegraphics[page=1, width=0.47\linewidth]{plan_mathe_lk11.pdf}}\hfill
	\fbox{\includegraphics[page=2, width=0.47\linewidth]{plan_mathe_lk11.pdf}}
\end{ausgabe}

Die \texttt{bemerkung} erscheint im Plan neben dem Termin, etwa für eine Stunde, die eine andere Lehrkraft übernimmt.



%%%%  KAPITEL 6  %%%%
\section{Mit einer KI eine neue Eingabedatei erstellen}

Weil die Eingabedatei kurz und streng gegliedert ist, kann eine KI sie gut nach Muster erstellen. Ein Beispiel-Prompt:

\begin{hinweis}[Beispiel-Prompt]
	\itshape Hier ist meine Datei plan\_informatik10.toml als Muster. Erstelle mir im exakt gleichen Format eine neue Datei für Mathematik Klasse 9 (zwei Klassen: 9a Mo und Do, 9c Di und Fr, jeweils wöchentlich eine Doppelstunde). Lernbereiche nach dem sächsischen Lehrplan mit Ustd., Themen kurz und allgemein, nach jedem Lernbereich eine Note. Sperre: Projektwoche 15.--19.03.2027.
\end{hinweis}

Titel und Unterrichtsstunden der Lernbereiche sollte man danach mit dem Lehrplan abgleichen.



%%%%  KAPITEL 7  %%%%
\section{Bitte prüfen}

\begin{itemize}
	\item \textbf{A/B-Zählung:} Voreingestellt zählen Ferienwochen nicht mit (\texttt{ab\_zaehlung = "schulwochen"}). Wechselt der A/B-Rhythmus an der eigenen Schule auch in den Ferien, auf \texttt{"kalenderwochen"} umstellen. Die Woche von \texttt{beginn} muss zu \texttt{woche\_bei\_beginn} passen.
	\item \textbf{Lernbereiche:} Titel und Ustd.\ in den Beispielen sind Platzhalter und müssen mit dem Lehrplan abgeglichen werden.
	\item \textbf{Bewegliche Ferientage:} Der 07.05.2027 (Freitag nach Himmelfahrt) ist in der Ferien-Datei auskommentiert. Ist er an der eigenen Schule beweglicher Ferientag, die drei Zeilen aktivieren.
	\item \textbf{Ferien:} Die Datei gilt für Sachsen im Schuljahr 2026/27 (Quelle: Sächsische Ferienverordnung, \url{https://www.revosax.sachsen.de}). Für andere Länder oder Jahre eine eigene Datei im selben Format anlegen.
\end{itemize}



%%%%  KAPITEL 8  %%%%
\section{Zusammenfassung}

\begin{itemize}
	\item Eine TOML-Datei beschreibt Klassen, Termine, Sperren und Lernbereiche; eine zweite die Ferien.
	\item \texttt{python3 stoffverteilungsplan.py plan.toml} schreibt den Plan, \texttt{pdflatex} übersetzt ihn.
	\item Die Bilanz im Terminal zeigt sofort, ob der Pflichtumfang passt.
	\item Ändert sich etwas (Klassenfahrt, Ausfall, neuer Stundenplan), ändert man eine Zeile und erzeugt den Plan neu.
\end{itemize}

\end{document}
